\section{The Tape Coding Conjecture}\label{sec:conjecture}

\subsection{Statement}

\begin{conjecture}[Tape Coding Conjecture]\label{conj:TCC}
\formal{Tape/DefsCE.lean}{RobustNoCode_TM}
For all $k$, $G\ge1$, $c$ and $C$, the following holds for all sufficiently large $n$ (with
$w=b(n)$, $L=n^2w$ and $r=\lceil nw/G\rceil$). No causal code solves complete exchange at the profile
capacities on the block network of a valid $k$-tape time--space diagram of length at most $cL$, for a
set of inputs of co-entropy at most $Cnw$.
\end{conjecture}

In words: on the block network of any linear-length tape diagram, complete exchange cannot be
realized by any coding scheme, even if the scheme only has to work for a $2^{-O(n\log n)}$ fraction
of the inputs. By \cref{thm:code-gap}, a solving code would beat multicommodity flow by a factor
$\Omega(\log n)$. The conjecture can therefore be read as ``on tape networks, coding beats routing by
$o(\log n)$'', and this reading is one of its equivalent forms (\cref{thm:tcc-forms}). The internal
name of the conjecture in the Lean library is \lean{RobustNoCode_TM} (``no code for \emph{complete
exchange} on \emph{tape-machine} networks, robust form'').

The \emph{plain} form, \lean{NoCode_TM}, asks the same only for codes that are correct on all inputs.
It refutes oblivious linear-time transposers (\cref{thm:oblivious}). Whether plain implies robust is
the subject of \cref{sec:plainrobust}.

\subsection{The reduction}

\begin{theorem}[The Tape Coding Conjecture refutes (T)]\label{thm:reduction}
\formal{Tape/Statements/Coding.lean}{not_linearTimeTransposable_of_robustNoCode, not_linearTimeTransposable_of_robustCE_TM, not_linearTimeTransposable_of_robustCE_TM_dist, codes_of_linearTimeTransposable, not_robustCE_TM_of_linearTimeTransposable}
If the Tape Coding Conjecture holds (in any of its robust forms of \cref{thm:tcc-forms}), then
$\neg\,$\lean{Transpose.LinearTimeTransposable}. In contrapositive form, and more strongly, a linear-time
transposer yields constants $k$, $G\ge1$, $c$, $C$ such that for \emph{every} large $n$ there is a valid
$k$-tape diagram of length at most $cL$ whose block network carries a causal code, within capacities in the
block profile, solving complete exchange on a set of co-entropy at most $Cnw$ (a refutation of the
conjecture needs such codes only for infinitely many $n$).
\end{theorem}

\begin{outline}
The formal proof goes through \cref{thm:coent-linear}. The conjecture rules out codes at co-entropy
$Cn\le Cnw$ for every $C$. A machine meeting \cref{def:T} with constant $c$ is correct within $2cL$ steps on every matrix input of every size $n\ge1$; after output buffering, the popular class of its coarse trace has co-entropy $O(n)$, and its
class code (\cref{lem:class-code}), with the state and the head positions carried by extra stationary
tapes (\cref{thm:fast-codes}(b)), solves complete exchange at the profile capacities at every large
size. This contradicts the conjecture at a single large $n$. The flow and distance forms are
equivalent (\cref{thm:tcc-forms}), and the second statement follows because $Cn\le Cnw$.
\end{outline}

Concretely, a proof \texttt{h} of \lean{RobustNoCode_TM} resolves~(T) by the file
\begin{center}
\texttt{theorem Transpose.answer : $\neg$ Transpose.LinearTimeTransposable :=}\\
\texttt{\ \ Transpose.Tape.not\_linearTimeTransposable\_of\_robustNoCode h}
\end{center}
which \leanfile{Verify.lean} (in \texttt{lean/}) accepts, once \leanfile{Transpose/Answer.lean} also imports
\leanfile{Transpose.Tape.Statements.Coding}. The reduction is one-way. A refutation of the conjecture is a
choice of constants $k$, $G$, $c$, $C$ with codes for infinitely many $n$; nothing makes these codes uniform,
finite-state, or local to slots, as a machine's code would be, and no converse of \cref{thm:reduction} is
proved. So a refutation would remove this coding obstruction but would not by itself decide (T).

\begin{theorem}[Consequences for every machine]\label{thm:tcc-consequences}
\formal{Tape/Statements/Coding.lean}{linearTime_fails_of_robustNoCode, no_fast_transposer_of_robustNoCode, no_fast_transposer_of_robustNoCode_bits, linearTime_fails_of_robustNoCode_logBits, not_linearTimeTransposable_of_robustNoCode_logBits}
Assume the Tape Coding Conjecture.
\begin{enumerate}[label=(\alph*)]
\item For every machine $M$ and constant $c$, for all large $n$ some $n\times n$ matrix is not
transposed within $c(L+1)$ steps: every machine fails at \emph{every} large size.
\item (Non-uniform, average case.) For all $k$, $g$, $c$, $C_0$, $C_1$ and all large $n$, no machine
with $k$ work tapes, $g+1$ work symbols and at most $2^{C_1nw}$ states (its transition table may
depend on $n$) halts within $cL$ steps with the transposed output on every input of a set of
co-entropy at most $C_0nw$.
\item (Extra bits.) For every $s$, if no code exists at the block capacities with $s(k,n)$ extra bits on
every memory edge (\lean{RobustNoCode_TM_bits}), then for all $k$, $g$, $c$, $C_0$ and all large $n$ no
machine with $k$ work tapes, $g+1$ work symbols and $\clog{(q+1)}+(k+1)\clog{(2cL+1)}\le s(k,n)$ halts
within $cL$ steps with the transposed output on every input of a set of co-entropy at most $C_0nw$.
For $s(k,n)=(3k+4)\clog n$ this hypothesis also implies (a) and $\neg$(T).
\end{enumerate}
\end{theorem}

\begin{outline}
Corollaries of the reduction at co-entropy $O(n)$ (\cref{thm:coent-linear}), whose single-size form
applies to each machine separately. The form with extra bits per memory edge is direct: a run's
payload on a memory edge is exactly its resident blocks, its state and its head positions.
\end{outline}

\subsection{Equivalent forms}

Two further forms make the connection with network coding explicit. The \emph{flow form}
\lean{RobustCE_TM} says: for every $\varepsilon>0$ and all large $n$, whenever a causal code solves
complete exchange within capacities $\mathrm{cap}$ in the block profile on a set of co-entropy at most
$Cnw$, the demand is routable at rate $1/(\varepsilon w)$ within $\mathrm{cap}$; that is, coding beats
multicommodity flow by a factor of at most $\varepsilon\log n$. The \emph{distance form}
\lean{RobustCE_TM_dist} replaces the flow rate by its volume bound:
$\sum_d\mathrm{rowDist}(d)\le\varepsilon w\sum_e\mathrm{cap}(e)$.

\begin{theorem}[Equivalent forms]\label{thm:tcc-forms}
\formal{Tape/Statements/Coding.lean}{robust_tfae, robustCE_TM_iff_robustNoCode, robustCE_TM_dist_iff_robustNoCode, robustNoCode_iff_attached, robustCE_TM_att_iff, robustCE_TM_att_iff_robustCE_TM, robustCE_TM_gamma_iff, robustCE_TM_gamma_per_iff, robustNoCode_of_gainBound, plain_tfae, ce_TM_iff_noCode, ce_TM_dist_iff_noCode, noCode_of_robustNoCode, ce_TM_of_robustCE_TM, ce_TM_dist_of_robustCE_TM_dist}
The following are equivalent: the Tape Coding Conjecture; its flow form; its distance form; the
no-code and flow forms restricted to diagrams whose input walk attaches every row; the no-code form
with the uniform profile of $3k$ blocks on every memory edge; and the no-code form with
$(3k+4)\clog n$ extra bits on every memory edge. The $\gamma$-form ``for some $\gamma(n)=o(\log n)$ with
$\gamma>0$ (one $\gamma$ for all constants, or one per choice of constants), every solving code's demand is
routable at rate $1/\gamma(n)$'' is equivalent as well. The explicit gain bound ``every code solving complete
exchange at the profile capacities has $E\sum_d\mathrm{rowDist}(d)\le w\sum_e\mathrm{cap}(e)$'', with $E$ as in
\cref{thm:code-gap}, implies the conjecture (the converse holds vacuously but is not formalized). The plain
forms (flow, distance, no code, for codes correct on all inputs) are equivalent to each other, and are
implied by the robust ones.
\end{theorem}

\begin{outline}
A solving code at the profile capacities makes the distance form fail by \cref{thm:code-gap}, and
weak duality (\cref{lem:volume}) turns any routing at rate $1/(\varepsilon w)$ into the distance form.
So each form is equivalent to ``no code''. The restriction to attached rows loses nothing: append a
final input sweep, which keeps every original node, edge and sink and the residency of every complete
slot, and let the code mask each node's input to its original rows (the time constant becomes $2c+3$).
The profile forms are \cref{thm:profile}.
\end{outline}

\begin{theorem}[The capacity profile is immaterial]\label{thm:profile}
\formal{Tape/Statements/Coding.lean}{robustNoCode_with_iff, robustNoCode_with_k_iff, exists_shift_iff_unbounded, robustNoCode_with_of_robustNoCode, robustNoCode_bits_iff, robustNoCode_logBits_iff, robustNoCode_iff_forall_cap, noCode_iff_forall_cap}
For every unbounded memory profile $m$ (equivalently, $3k\le m(k+a)$ for some $a$, for every $k$; for
instance $m(k)=3k$, $3k+1$, or one block per tape), the no-code statement with $m(k)$ blocks on every
memory edge is equivalent to the Tape Coding Conjecture. So is the no-code statement with $s(k,n)$
extra bits on every memory edge, for every $s$ with $s(k,n)\le C(k)\,nw$; for instance
$s=(3k+4)\clog n$. The conjecture stated at the profile capacities is equivalent to ruling out codes
within every capacity in the profile, and the same holds for the plain form.
\end{theorem}

\begin{outline}
Extra bits and extra blocks are traded for stationary tapes (\cref{thm:fast-codes}(b)): a stationary
tape changes no node, edge, source, sink or chord, and its resident block can carry any payload of at
most one block. Conversely the per-node profile is at most any uniform profile with enough tapes.
\end{outline}

\subsection{Weaker sufficient conditions}

The conjecture asks for no codes at co-entropy $Cnw$, at all large $n$. Much less suffices to
refute~(T).

\begin{theorem}[Co-entropy $O(n)$ suffices]\label{thm:coent-linear}
\formal{Tape/Statements/Coding.lean}{bufTM_run, codes_coentLinear_of_linearTimeTransposable, not_linearTimeTransposable_of_coentLinear, linearTime_fails_of_coentLinear, no_fast_transposer_of_coentLinear, not_linearTimeTransposable_of_coentLinear_end, not_linearTimeTransposable_of_coentLinear_end_io, coentLinear_of_robustNoCode, noCode_of_coentLinear}
Write $\mathrm{NoCode}(m)$ for the statement ``for all $k$, $G\ge1$, $c$ and all large $n$, no causal
code solves complete exchange at the profile capacities on a valid $k$-tape diagram of length at most
$cL$, for a set of co-entropy at most $m(n)$'' (\lean{NoCode_TM_coent}).
\begin{enumerate}[label=(\alph*)]
\item (T) yields codes at co-entropy $Cn$ (not $Cnw$), on diagrams that append their output in their
last $L$ steps.
\item Hence, if $\mathrm{NoCode}(Cn)$ holds for every $C$, then $\neg$(T); every machine fails at every
large size; and the non-uniform average-case statement of \cref{thm:tcc-consequences}(b) holds with
co-entropy $C_0n$. It suffices to rule out codes on diagrams with the output at the end, and for
infinitely many $n$.
\item This hypothesis lies between the plain and the robust form: the Tape Coding Conjecture implies
it, and it implies $\mathrm{NoCode}(0)$, the plain form.
\end{enumerate}
\end{theorem}

\begin{outline}
First make the machine write its output only at the end (\lean{Impl.bufTM}): it writes the output on an
extra tape, rewinds, and copies it out, halting at $T+2L+2$ with the same output. Then the output length
at every slot is fixed by the halting time, and a \emph{coarse trace} suffices: block and row moves in
$\{-1,0,1\}$ together with the resident-block and touched-row patterns within $\pm1$, at most
$24^{k+1}$ values per slot by (V1) and (V2) of \cref{lem:residency} (for the input head: a slot is at most
one row long), which determine the block-level trace. The popular coarse
class therefore has co-entropy $O(n)$, and the class code applies unchanged.
\end{outline}

\begin{theorem}[One superlinear level, infinitely many sizes]\label{thm:io}
\formal{Tape/Statements/Coding.lean}{not_linearTimeTransposable_of_coent, not_linearTimeTransposable_of_coent_nw, not_linearTimeTransposable_of_coentLinear_io, not_linearTimeTransposable_of_coent_io, coent_io_of_coent, coentLinear_io_of_robustNoCode_io, coentLinear_io_iff_exists_coent_io, coent_io_nw_of_robustNoCode_io, not_linearTimeTransposable_of_robustNoCode_io, not_linearTimeTransposable_of_robustCE_TM_io, robustNoCode_io_of_robustNoCode, robust_io_tfae, no_fast_transposer_of_coent, traceClass_coent_lt}
If $f(n)\ge Cn$ eventually for every $C$ (for instance $f(n)=nw$), then $\mathrm{NoCode}(f)$ refutes
(T); so does $\mathrm{NoCode}(f)$ for infinitely many $n$. The infinitely-often forms of the conjecture
(flow, distance, no code) are equivalent to each other, formally weaker than the conjecture, and still
refute (T). For infinitely many $n$, the single superlinear level and the linear family are equivalent:
$\mathrm{NoCode}(Cn)$ holds for infinitely many $n$, for every $C$, if and only if $\mathrm{NoCode}(f)$
holds for infinitely many $n$ for some $f$ with $f(n)\ge Cn$ eventually for every $C$. Under
$\mathrm{NoCode}(m)$, for all large $n$: no machine with $k$ work tapes, $g+1$ symbols and $2^{O(nw)}$
states is correct within $cL$ steps on a set of co-entropy $m_0$ with $m_0+Kn\le m(n)$, for a constant $K$
depending only on $k$, $g$ and $c$; and every trace class on which such a machine is correct within $cL$
steps has co-entropy larger than $m(n)$.
\end{theorem}

\begin{outline}
(T) gives codes at every large $n$, so a single failing size per choice of constants contradicts it.
For a single level, $Cn\le f(n)$ from some point on. Of the equivalence of the single-level and the
linear infinitely-often families, one direction is monotonicity in the level and the other a diagonal
argument.
\end{outline}

\begin{theorem}[Further equivalent and sufficient forms]\label{thm:rungs}
\formal{Tape/Statements/Coding.lean}{coentLinear_of_robustCE_TM_lin, robustCE_TM_lin_iff, not_linearTimeTransposable_of_robustCE_TM_lin, hintCodes_linear_of_linearTimeTransposable, hintLinear_iff_coentLinear, not_linearTimeTransposable_of_noCode_hintLinear}
The flow form at co-entropy $Cn$ is equivalent to $\mathrm{NoCode}(Cn)$ for every $C$, and refutes (T).
(T) gives codes indexed by an $O(n)$-bit oracle hint (\cref{sec:plainrobust}), and ruling out codes with
a $Cn$-bit hint for every $C$ is equivalent to $\mathrm{NoCode}(Cn)$ for every $C$, so it refutes (T)
as well.
\end{theorem}

\begin{outline}
As in \cref{thm:tcc-forms}, at the linear rung; hints and co-entropy are interchanged by
\cref{thm:hint}.
\end{outline}

\subsection{Oblivious machines}

\begin{theorem}[The plain conjecture refutes oblivious (T)]\label{thm:oblivious}
\formal{Tape/Statements/Coding.lean}{codes_of_obliviousLinearTimeTransposable, not_obliviousLinearTimeTransposable_of_noCode, not_obliviousLinearTimeTransposable_of_ce_TM, not_obliviousLinearTimeTransposable_of_ce_TM_dist, not_obliviousLinearTimeTransposable_of_noCode_io, noCode_io_of_robustNoCode_io}
Let \lean{ObliviousLinearTimeTransposable} be (T) for a machine that is oblivious at every size. An
oblivious linear-time transposer gives codes for independent uniform sources (all inputs). Hence the
plain form of the conjecture, in any of its forms, refutes oblivious (T); so does the plain no-code
form for infinitely many $n$. The robust infinitely-often form implies the plain infinitely-often form.
\end{theorem}

\begin{outline}
For an oblivious machine all inputs form one trace class with $D=0$ (\cref{lem:trace-obl}), so the
class code works on all inputs, with the state and the head positions on one extra stationary tape
(\cref{thm:fast-codes}(b)).
\end{outline}
